\section{The spectral side of the trace formula}

The spectral side of the trace formula comprises the collection of terms
associated with the representations of the group. This side, though
analogous to the already discussed geometric side of the trace formula in
terms of the overall structure, requires more background to state.
Essentially, we will have to leave the entire theory of Eisenstein series as
``known'' to the reader. Fortunately, there are several excellent
expositions of this theory in the case of $SL(2)$ and related groups:
[L1], [K2], [KP] for example.

A general remark may be in order here. The spectral theory used in [A2]
relies on the general theory of Eisenstein series on a quotient $\Gamma
\backslash H$ developed in [L1], where $H$ denotes a Lie group with finitely
many connected components and $\Gamma $ denotes a discrete subgroup
satisfying certain assumptions. By passing to the archimedian place, the
quotient $G({\mathbb{Q}})\backslash \overline{G({\mathbb{A}})^1}$ may be regarded as a
finite union of quotients $\Gamma _i\backslash H$ as in [L1]
(see Lemma 1.7 above). Thus we may apply the results of [L1] here.

We begin with a brief overview of some of the background. Some results we
will eventually need are the decomposition (4.5), Proposition 4.14
and the definition following it, and Theorem 6.15. Other than (4.5),
much of 4.1.1 can be ignored or skimmed over as we sketch the
(slightly different) decompositions used in [A1] and in [A2,A3].
We only need the decomposition (over a set ${\cal{X}}$ introduced below)
used in [A2,A3].

\subsection{The spectral decomposition}

A basic result in representation theory [Ro, ch. 14], [L2, Lemma 3.2] is the
following

\begin{lemma}
Let $\Gamma \subset G$\ be a discrete subgroup of a locally compact
unimodular group $G$. The right regular representation of $G$
acting on $L^2(\Gamma \backslash G)$
decomposes as a (not necessarily discrete)
Hilbert sum of irreducible representations, each occuring with finite
multiplicity.
\end{lemma}

For $H<\overline{G({\mathbb{A}})}$ a locally
compact totally disconnected group of
type 1, define $\Pi_u(\overline{H})$ as in section 3.2 above.

Another way to look at the spectral decomposition
is adelically: The structure of the
unitary dual $\Pi_u(\overline{G({\mathbb{A}})})$ can be determined from the
structure of $\Pi_u(\overline{G({\mathbb{Q}}_p)})$.
(The proof of this is given in
[Ge2] and is a consequence of a theorem of Gelbart
mentioned in Example 5.3 below.) The {\bf
spectrum} of $G({\mathbb{Q}} )\backslash \overline{G({\mathbb{A}})}$ is the set of
elements of $\Pi_u(\overline{G({\mathbb{A}})})$
which occur in the representation
space $L^2(G({{\mathbb{Q}}} )\backslash \overline{G({\mathbb{A}})})$, on which
$\overline{G({\mathbb{A}})}$ acts on the right by translation.

\subsubsection{The partial decomposition in the adelic case: cuspidals}

For each parabolic $P=MN$ of $G$ defined over ${\mathbb{Q}}$, define 
\begin{equation}
\begin{array}{c}
L^2(M({\mathbb{Q}})\backslash \overline{M({\mathbb{A}})^1})_{cusp} \\ 
=\{\phi \in L^2(M({{\mathbb{Q}}})\backslash \overline{M({\mathbb{A}})^1})
\ {\rm \ genuine\ }|\;{\rm for\;any\;parabolic}\;
P_1\subset P,\;P_1\neq P, \\ 
{\rm we\;have}
\int_{N_1({\mathbb{Q}})\cap M({\mathbb{Q}})\backslash
N_1({\mathbb{A}})\cap M({\mathbb{A}})}\phi (nm)dn=0, \\ 
{\rm for\;almost\;all}\;m\in \overline{M({\mathbb{A}})^1}\}.
\end{array}
\end{equation}
This is the space of genuine {\bf cusp forms} on
$\overline{M({\mathbb{A}})^1}$.

\begin{example}
(representations of adele groups) {\sf Case} $P=M=G=SL(2)$:
The genuine unitary representations $\pi $ of
$\overline{G({{\mathbb{A}}})}$ are all
``factorizable'', i.e.,
$\pi =\otimes _{p\leq \infty}\pi _p$ on
$\overline{G({{\mathbb{A}}})}$, where
$\pi _p$ is a genuine irreducible unitary representation of
$\overline{G({\mathbb{Q}}_p)}$ and almost all $\pi _p$
are unramified (i.e., has a non-trivial
$\overline{K_p}$-fixed vector). Conversely, given a collection 
\[
\{\pi _p\;|\;\pi _p\in \Pi _u(\overline{G({\mathbb{Q}}_p)}),
\;{\rm almost\;all}\;\pi _p\;{\rm are\;unramified}\},
\]
the restricted tensor product
$\pi =\otimes _{p\leq \infty}\pi _p$
defines an element of $\Pi _u(\overline{G({\mathbb{A}})})$. For a proof,
see for example, [Ge2, \S 3].

\end{example}

\begin{proposition}
For each $\rho \in \Pi _u(\overline{M({\mathbb{A}})^1})$
there is an integer $m(\rho )\geq 0$ such that 
\[
L^2(M({{\mathbb{Q}}})\backslash \overline{M({\mathbb{A}})^1})_{cusp}\cong
\bigoplus_{\rho \in \Pi _u(\overline{M({\mathbb{A}})^1})}m(\rho )\rho .
\]
\end{proposition}

If
$\rho \in \Pi _u(\overline{M({\mathbb{A}})^1})$ and if $m(\rho )\neq 0$
then we call $\rho $ a genuine {\bf cuspidal
representation} of $\overline{M({\mathbb{A}})^1}$.

Two parabolics $P=MN$ and $P^{\prime }=M^{\prime }N^{\prime }$ are said to
be {\bf associate} parabolics if 
\[
\Omega ({\mathfrak{a}},{\mathfrak{a}}^{\prime })
=\{s\in \Omega \;|\;s{\mathfrak{a}}_{{\mathbb C}}
={\mathfrak{a}}_{{\mathbb C}}^{\prime }\}\neq \emptyset . 
\]
Here ${\mathfrak{a}}_{\mathbb{C}}$ denotes the
complexification of the Lie algebra
${\mathfrak{a}}$ of the center $A$ of $M$. In particular, a parabolic is
always associate to its opposite.
We identify the Weyl group $\Omega$ with a
subgroup of $G({\mathbb{Q}})$ (hence as a subgroup of
$\overline{G({\mathbb{A}})}$) by $s\longmapsto w_s\in G({\mathbb{Q}})$.
If $P,P^{\prime }$ are associate and if 
$\rho \in \Pi _u(\overline{M({\mathbb{A}})^1}),
\;\rho ^{\prime }\in \Pi _u(\overline{M^{\prime }({\mathbb{A}})^1})$
are cuspidal then we call $(M,\rho )$ {\bf equivalent} to
$(M^{\prime },\rho ^{\prime })$ if there exists an
$s\in \Omega ({\mathfrak{a}}, {\mathfrak{a}}^{\prime })$ such that
$s\rho $ is unitarily equivalent (as a
representation of $\overline{M^{\prime }({\mathbb{A}})^1})$
to $\rho ^{\prime }$, where 
\[
(s\rho )(m^{\prime })=\rho (w_s^{-1}m^{\prime }w_s),
\;\;\;\;m^{\prime }\in \overline{M^{\prime }({{\mathbb{A}}})^1}. 
\]
{\bf Notation}:
Let ${\cal{X}}={\cal{X}}^G$ denote the set of equivalent classes of such
pairs $(M,\rho )$.
For each $\chi \in {\cal X}$, we denote by
${\cal P}_\chi $ the corresponding class of associated parabolics.
Let ${\cal{X}}(G)$ denote the set of equivalent classes of
cuspidal pairs $(G,\rho )$.

\begin{example}
Two parabolics $P,P^{\prime }$ of $G=SL(2)$ with Levi $A$
are associate if and only they are opposite.
Therefore $(M,\rho )$ is equivalent to $(M^{\prime },\rho ^{\prime })$
if and only if $M=M^{\prime }=G$ and $\rho =\rho ^{\prime }$
(as equivalence classes of representations of
$\overline{G({\mathbb{A}})}$),
or $M=M'=A$ and $\rho ={\rho'}^w$, where 
$w=
\left(
\begin{array}{cc}
0 &-1\\
1&0
\end{array}
\right)$
denotes the non-trivial element of the Weyl group of $A$.
\end{example}

For each $P=MN$ and each $\chi \in {\cal X}$, we let 
\begin{equation}
L^2(M({{\mathbb{Q}}})\backslash \overline{M({\mathbb{A}})^1})_{cusp,\chi }
\cong
\bigoplus_{\rho \in \Pi _u(\overline{M({\mathbb{A}})^1}) ,
(M,\rho )\in \chi }
m(\rho )\rho .
\end{equation}
Note that if $P\notin {\cal P}_\chi $ then this sum is empty (which, by
convention, we define to be $0$).

The above theorem has the following

\begin{corollary}
$L^2(M({\mathbb{Q}})\backslash \overline{M({\mathbb{A}})^1})_{cusp}
=\bigoplus_{\chi \in {\cal X}}
L^2(M({\mathbb{Q}})\backslash \overline{M({\mathbb{A}})^1})_{cusp,\chi }$.
In particular, 
\[
L^2(G({\mathbb{Q}})\backslash \overline{G({\mathbb{A}})})_{cusp}
=\bigoplus_{\chi \in {\cal X}}
L^2(G({\mathbb{Q}})\backslash \overline{G({\mathbb{A}})})_{cusp,\chi }.
\]
\end{corollary}

In other words, the space of cuspidal automorphic functions decomposes into
``irreducible'' subspaces on which the action of $\overline{G({\mathbb{A}})}$ is the
more-or-less the same.

\subsubsection{The partial decomposition in the adelic case: theta functions}

We seek an analogous result for all of
$L^2(G({\mathbb{Q}})\backslash \overline{G({\mathbb{A}})})_{gen}$.

Let ${\mathfrak{Z}}_M$ denote the center of the
universal enveloping algebra ${\mathfrak{U}}_M$ of $M$. Each
$X\in {\mathfrak{Z}}_M$ defines a densely defined
unbounded operator on
$L^2(M({\mathbb{Q}})\backslash \overline{M({\mathbb{A}})^1})_{cusp}$.

Let ${\cal H}_P^0$ denote 
\begin{equation}
\begin{array}{c}
{\cal H}_P^0
=\{\varphi :N({\mathbb{A}})M({\mathbb{Q}})
\backslash \overline{G({\mathbb{A}})}\rightarrow {\mathbb{C}}
\ {\rm genuine}\;|\; \\ 
(i)\;\forall x\in \overline{G({\mathbb{A}})},
\;{\rm the\;function}\;m\longmapsto \varphi (mx),
\;m\in M({\mathbb{A}}), \\ 
{\rm is}\;{\mathfrak{Z}}_{M({\mathbb{R}})}{\rm -finite}, \\ 
(ii)\;{\rm the\;span\;of\;all\;functions}
\;\varphi _k:x\longmapsto \varphi (xk),\;k\in \overline{K}, \\ 
{\rm is\;finite\;}{\rm dimensional}, \\ 
(iii)\;||\varphi ||^2=
\int_{\overline{K}}
\int_{M({\mathbb{Q}})\backslash \overline{M({\mathbb{A}})^1}}
|\varphi (mk)|^2dmdk<\infty \} .
\end{array}
\end{equation}
Let ${\cal H}_P$ denote the closure of ${\cal H}_P^0$ with respect to this
norm in (iii).

For $\varphi \in {\cal H}_P$,
$\Lambda \in {\mathfrak{a}}_{{\mathbb C}},\;x,y\in \overline{G({\mathbb{A}})}$,
define $I_P(\Lambda ):\overline{G({\mathbb{A}})}\rightarrow End({\cal H}_P)$ by 
\[
(I_P(\Lambda ,y)\varphi )(x)=\varphi (xy)\exp (<\Lambda +\rho
_P,H_P(xy)>)\exp (-<\Lambda +\rho _P,H_P(x)>). 
\]
Here $\rho _P\in {\mathfrak{a}}$ is the element such that
$|\det (Ad(m))_{{\mathfrak{n}}({\mathbb A})}|
=\exp (<2\rho _P,H_P(m)>),\;m\in \overline{M({\mathbb{A}})}$.

For each $P=MN$ and each $\chi \in {\cal X}$, we define
$L^2(M({\mathbb{Q}})\backslash \overline{M({\mathbb{A}})^1})_\chi $
as follows. If there is no $P_1\in {\cal P}_\chi $
with $P_1\subset P$ then we define
$L^2(M({{\mathbb{Q}}})\backslash \overline{M({\mathbb{A}})^1})_\chi $
to be $0$. Otherwise, let $\Psi (P_1,P,\chi )$ denote the space of
all genuine complex-valued functions
$\psi :N({\mathbb{A}})M_1({\mathbb{Q}})\backslash \overline{G({\mathbb{A}})}
\rightarrow {\mathbb C}$, $P_1=N_1M_1\in {\cal{P}}_\chi $ 
with $P_1\subset P$, such that

\begin{enumerate}
\item  $\psi (amk)$, as a function of 
$a\in A_1({\mathbb{Q}})\backslash \overline{A_1({\mathbb{A}})}$ (fixed 
$m\in M_1({\mathbb{Q}})\backslash \overline{M_1({\mathbb{A}})^1},\;k\in K$),
is compactly supported,

\item  $\psi (amk)$, as a function of $k\in \overline{K_\infty}$,
belongs to the
vector space defined by the span of the matrix coefficients of some
irreducible representation $W$ of $\overline{K_\infty}$,

\item  $\psi (amk)$, as a function of $m\in M$, belongs to 
$L^2(M({{\mathbb{Q}}})\backslash \overline{M({\mathbb{A}})^1})_{cusp,\chi }$.
\end{enumerate}

Given $\psi \in \Psi (P_1,P,\chi )$, $P_1\subset P$, let 
\[
\psi ^{\wedge ,M}(m)=
\sum_{\delta \in P_1({\mathbb{Q}})\cap
M({\mathbb{Q}})\backslash M({\mathbb{Q}})}
\psi (\delta m),\;\;\;\;\;
m\in M({\mathbb{Q}})\backslash \overline{M({\mathbb{A}})^1}, 
\]
so that
$\psi ^{\wedge ,M}\in
L^2(M({\mathbb{Q}})\backslash \overline{M({\mathbb{A}})^1})$
[L1,\ Lemma 1]. Let
$L^2(M({\mathbb{Q}})\backslash \overline{M({\mathbb{A}})^1})_\chi $ denote
the $L^2$-closure of the span of all the
$\psi ^{\wedge ,M}$ with $\psi \in \Psi (P_1,P,\chi )$,
some $P_1\subset P$ as above. Lemma 2 of [L1] implies 
\begin{equation}
L^2(M({\mathbb{Q}})\backslash \overline{M({\mathbb{A}})^1})
=\bigoplus_{\chi \in {\cal X}^G}
L^2(M({\mathbb{Q}})\backslash \overline{M({\mathbb{A}})^1})_\chi .
\end{equation}

\subsection{Eisenstein series in the adelic case}

For $\varphi \in {\cal H}_P^0,\;\lambda \in {\mathfrak{a}}_{{\mathbb C}},\;s\in
\Omega ({\mathfrak{a}},{\mathfrak{a}}_1),\;x\in \overline{G({\mathbb{A}})}$, let
$M(s,\lambda)$ denote the intertwining integral
\begin{equation}
\begin{array}{c}
M(s,\lambda )\phi(x) \\ 
=\exp (-<s\lambda +\rho _1,H_P(x)>)\times \;\;\;\;\;\;\;\;\;\;\;\;\; \\ 
\;\;\;\;\;\;\times
\int_{N_1({\mathbb{A}})\cap w_sN({\mathbb{A}})w_s^{-1}\backslash
N_1({\mathbb{A}})}\varphi (w_s^{-1}nx)
\exp (<\lambda +\rho_P,H_P(w_s^{-1}nx)>)dn,
\end{array}
\end{equation}
with a suitably normalized measure $dn$.
We define $M(s,\pi_\zeta)$ below (see (5.10)).

Let ${\mathfrak{a}}^{+}$ denote the positive Weyl chamber ${\mathfrak{a}}^{+}=\{X\in 
{\mathfrak{a}}\;|\;\alpha (X)>0,\;\forall \alpha >0\}$.
For $\varphi \in {\cal{H}}_P^0,
\;\lambda \in {\mathfrak{a}}_{{\mathbb C}},
\;Re(\lambda )\in \rho _P+{\mathfrak{a}}^{+},
\;x\in \overline{G({\mathbb{A}})}$, let 
\begin{equation}
E(x,\varphi ,\Lambda )=
\sum_{\delta \in P({\mathbb{Q}})\backslash G({\mathbb{Q}})}
\varphi (\delta x)\exp (<\lambda +\rho _P,H_P(\delta x)>).
\end{equation}
This is the {\bf Eisenstein series} associated to $P$ and $\varphi $ and is
known to converge absolutely. 

Let $P=MN$ denote a parabolic and let
$\pi \in \Pi _u(\overline{M({\mathbb{A}})^1})$.
Following [A2], we define 
\begin{equation}
\begin{array}{c}
{\cal H}_P^0(\pi )=\{\varphi :N({\mathbb{A}})M({\mathbb{Q}})
\backslash \overline{G({\mathbb{A}})}\rightarrow {\mathbb C}\;
{\rm genuine\ }|\; \\ 
(i)\;\phi \;{\rm is\;right}\;K{\rm -finite,} \\ 
(ii)\;\forall x\in \overline{G({\mathbb{A}})},
\;{\rm the\;function}\;m\longmapsto \varphi (mx)
\;(m\in \overline{M({{\mathbb{A}}})^1}) \\ 
{\rm is\;a\;matrix\;coefficient\;of}\;\pi , \\ 
(iii)\;||\varphi ||^2
=\int_K\int_{M({\mathbb{Q}})\backslash \overline{M({\mathbb{A}})^1}}|
\varphi (mk)|^2dmdk<\infty \}.
\end{array}
\end{equation}

Let ${\cal H}_P(\pi )$ denote the completion of ${\cal H}_P^0(\pi )$ with
respect to the norm $||...||$ in (iii). For
$\phi \in {\cal H}_P(\pi )$, $\varsigma \in i{\mathfrak{a}}^{*},$ let
\begin{eqnarray*}
\phi _\varsigma (x)
&=&\phi (x)e^{\varsigma (H_P(x))},\;\;\;x\in \overline{G({\mathbb{A}})},
\\
\pi _\varsigma  &=&\pi \otimes \chi _\varsigma ,\;\;\;\chi _\varsigma
(x)=e^{\varsigma (H_P(x))},\;\;x\in \overline{G({\mathbb{A}})}, \\
(I_P(\pi _\varsigma ,y)\phi _\varsigma )(x) &=&\phi _\varsigma (xy)\delta
_P(xy)^{1/2}\delta _P(x)^{-1/2},\;\;\;x,y\in \overline{G({\mathbb{A}})}.
\end{eqnarray*}

Let $P=MN$ denote a parabolic, let $\pi \in \Pi _u(\overline{M({\mathbb{A}})})$,
let $f$ be a function on $\overline{M({\mathbb{A}})}$,
and let $\zeta\in {\mathfrak{a}}^*_{\mathbb{C}}$. Define $\pi_\zeta$ by
\begin{equation}
\begin{array}{c}
f_\zeta (m)=f(m)e^{\zeta(H_M(m))},\\
\pi_\zeta(m)=\pi(m)e^{\zeta(H_M(m))},\ \ \ m\in \overline{M({\mathbb{A}})}.
\end{array}
\end{equation}

The space of (genuine) matrix coefficients of $\pi$, which we denote by
${\cal A}(\pi)$, is isomorphic to that of $\pi_\zeta$, for
$\zeta\in i{\mathfrak{a}}^*$:
$$
{\cal A}(\pi)\cong {\cal A}(\pi_\zeta).
$$
(proof: We may let $\pi,\ \pi_\zeta$ have the same representation space
$V$ with inner product $(...,...)_V$. If $v_1,v_2\in V$ and
$f(g)=(\pi(g)v_1,v_2)_V$ then $f_\zeta\in {\cal A}(\pi_\zeta)$ and every
element of ${\cal A}(\pi_\zeta)$ arises in this way in a unique
fashion. $\Box$) We have a natural isomorphism
$$
\begin{array}{c}
{\cal H}_P(\pi)\rightarrow {\cal H}_P(\pi_\zeta),\\
\phi\longmapsto \phi_\zeta,
\end{array}
$$
For $\pi \in \Pi _u(\overline{M({\mathbb{A}})})$,
$\phi \in {\cal H}_P^0(\pi )$,
$\varsigma \in {\mathfrak{a}}_{P,{\mathbb C}}^{*}$,
$s\in \Omega ({\mathfrak{a}},{\mathfrak{a}}^{\prime})$,
$Re(\varsigma )\in \rho _P+({\mathfrak{a}}_P^{*})^{+}$, define
\begin{equation}
\begin{array}{c}
E(x,\phi _\varsigma )
=\sum_{\delta \in P({\mathbb{Q}})\backslash G({\mathbb{Q}})}
\phi _\varsigma (\delta x)\delta _P(\delta x)^{1/2}, \\
(M(s,\pi _\varsigma )\phi _\varsigma )(x)
=\delta_P(x)^{-1/2}
\int_{N^{\prime }({\mathbb{A}})\cap w_sN({\mathbb{A}})w_s^{-1}\backslash
N^{\prime }({\mathbb{A}})}
\phi _\varsigma (w_s^{-1}nx)\delta_P(w_s^{-1}nx)^{1/2}dn.
\end{array}
\end{equation}
Analogously, if $P\subset P_2$, $\varsigma \in {\mathfrak{a}}_{{\mathbb C}}^{*}$ is
sufficiently large then define
\[
E_{P_2}(x,\phi _\varsigma )
=\sum_{\delta \in P({\mathbb{Q}})\backslash P_2({\mathbb{Q}})}
\phi _\varsigma (\delta x)\delta _P(\delta x)^{1/2}.
\]

Let
\[
\begin{array}{c}
{\cal H}_P(\pi )_\chi
=\{\phi \in {\cal H}_P(\pi )\;|\;m\longmapsto \varphi (mx)
\;(m\in M({\mathbb{Q}})\backslash \overline{M({\mathbb{A}})}^1) \\
{\rm \;\;\;\;\;\;\;\;\;\;\;\;belongs\;to}
\;L^2(M({\mathbb{Q}})\backslash \overline{M({\mathbb{A}})^1})_\chi \}.
\end{array}
\]

\begin{proposition}
${\cal H}_P(\pi )=\bigoplus_{\chi \in {\cal X}}{\cal H}_P(\pi )_\chi $.
\end{proposition}

Let $K\subset G({\mathbb{A}}_f)\cap K_0$ be an open compact subgroup and let 
\begin{equation}
{\cal H}_P(\pi )_{\chi ,\overline{K}}\subset {\cal H}_P(\pi )_\chi 
\end{equation}
denote the subspace which is right invariant under
$\overline{K}$. Let $W$ denote a genuine
irreducible representation of $\overline{K_\infty} $ and let 
\begin{equation}
{\cal H}_P(\pi )_{\chi ,\overline{K},W}
\subset {\cal H}_P(\pi )_{\chi ,\overline{K}}
\end{equation}
denote the subspace of functions $\phi $ such that
$k\longmapsto \phi (xk)\;\;(k\in \overline{K_\infty} )$
is a matrix coefficient of $W$.

\begin{proposition}
$\dim {\cal H}_P(\pi )_{\chi ,\overline{K},W}<\infty .$
\end{proposition}

We apply this to fixing a choice of a basis:
For each connected component $\Pi^0\subset
\Pi_u(\overline{M({\mathbb{A}})}^1)$, fix an element $\pi^0\in \Pi^0$.
Every other $\pi\in\Pi^0$ may be uniquely written as
$\pi=\pi^0_\zeta$, for some $\zeta\in i{\mathfrak{a}}^*_P$.
Let ${\cal B}_P(\pi^0 )_\chi $
denote a basis of ${\cal H}_P(\pi^0 )_\chi $ satisfying the following
properties:

\begin{enumerate}
\item  each $\phi \in {\cal B}_P(\pi^0 )_\chi $ belongs to
${\cal H}_P(\pi^0)_{\chi ,\overline{K},W}$
for some $K,W$ (possibly depending on $\phi $),

\item  the elements of ${\cal B}_P(\pi^0 )_\chi $
form an orthonormal set with
respect to the inner product
\[
(\phi _1,\phi _2)=
\int_K\int_{M({\mathbb{Q}})\backslash \overline{M({\mathbb{A}})}^1}
\phi _1(mk)\overline{\phi _2(mk)}dmdk,
\]

\item  if $\varsigma \in i{\mathfrak{a}}^{*}$ then 
\[
{\cal B}_P(\pi^0_\varsigma )_\chi =\{\phi _\varsigma \;|\;
\phi \in {\cal B}_P(\pi^0 )_\chi \}.
\]
\end{enumerate}

${\cal B}_P(\pi)_\chi$ is a basis for
${\cal H}_P(\pi)_\chi$, if $\pi=\pi^0$, and is a subset of
${\cal H}_P(\pi)$ in general.

Let ${\cal P}$ denote a class of associated parabolics. Let
$L_{{\cal P}}^{\wedge }$ denote the
set of collections $F=\{F_P\;|\;P\in {\cal P}\}$ of
measurable functions
$F_P:\Pi _u(\overline{M({\mathbb{A}})})\rightarrow {\cal H}_P$ such
that

(i) if $s\in \Omega ({\mathfrak{a}},{\mathfrak{a}}_1)$ then 
\[
F_{P_1}(s\pi )=M(s,\pi )F_P(\pi ),
\]

(ii) if $N(A)$ denotes the number of chambers in ${\mathfrak{a}}$
(which is equal to 1 for $SL(2)$) then 
\[
||F||^2=\sum_{P\in {\cal{P}}}{1\over 2\pi }^{dim(A)}
\int_{\Pi _u(\overline{M({\mathbb{A}})})}
||F_P(\pi )||^2d\pi <\infty .
\]

The following result describes some basic properties of Eisenstein series.

\begin{theorem} For
$s\in \Omega ({\mathfrak{a}},{\mathfrak{a}}_1)$, we have
{\bf \ }(Langlands [A2, p. 927])

(a)\ If $\varphi \in {\cal H}_P^0(\pi )$\ then
$E(x,\varphi _\varsigma )$ and
$M(s,\pi _\varsigma )\varphi _\varsigma $
can be analytically continued
to a meromorphic function on ${\mathfrak{a}}_{{\mathbb C}}^{*}$. On 
$i{\mathfrak{a}}$ the
Eisenstein series $E(x,\varphi_\varsigma )$
is regular and the intertwining operator $M(s,\pi_\varsigma )$
is unitary. For $\overline{K_0}$-finite
$f\in C_c^\infty (\overline{G({\mathbb{A}})})_{gen}$,
$t\in \Omega ({\mathfrak{a}}_1,{\mathfrak{a}}_2)$, we have 
\[
\begin{array}{c}
(i)\;\;\;\;E(x,I_P(\pi ,f)\varphi )
=\int_{\overline{G({\mathbb{A}})}}f(y)E(xy,\varphi )dy,
\\ 
(ii)\;\;\;\;\;E(x,M(s,\pi )\varphi )=E(x,\varphi ), \\ 
(iii)\;\;\;\;M(ts,\pi )=M(t,s\pi )M(s,\pi ).
\end{array}
\]
(b) There is a dense subspace of $L_{{\cal P}}^{\wedge }$\ on which the map 
\[
F\longmapsto \sum_{P\in {\cal P}}n(A)^{-1}
\left( \frac 1{2\pi i}\right)^{\dim (A)}
\int_{\Pi _u(\overline{M({\mathbb{A}})})}E(x,F_P(\pi ))d\pi 
\]
is defined. This extends to all of $L_{{\cal P}}^{\wedge }$\ and defines a
unitary map of $L_{{\cal P}}^{\wedge }$\ onto a closed subspace 
\[
L_{{\cal P}}^2(G({{\mathbb{Q}}})\backslash \overline{G({\mathbb{A}})})
\subset L^2(G({\mathbb{Q}})\backslash \overline{G({\mathbb{A}})})_{gen}.
\]
Moreover, 
\[
L^2(G({\mathbb{Q}})\backslash \overline{G({\mathbb{A}})})_{gen}
=\bigoplus_{\cal{P}}
L_{\cal{P}}^2(G({\mathbb{Q}})\backslash \overline{G({\mathbb{A}})}).
\]
\end{theorem}

